
\documentclass[11pt]{article}
\title{HW1, Ghafouralhosseini}
\author{سیده فاطمه غفورالحسینی}
\date{October 2026}

\usepackage[utf8]{inputenc}
\usepackage[margin=1in]{geometry}
\usepackage{amsmath,amssymb}

\newcommand{\dd}{\mathrm{d}}

\begin{document}

\begin{center}
{\Large\bfseries Lattice propagator and the continuum limit}
\end{center}
Fatemeh Ghafouralhosseini, 404210037
\section*{Setup}

Consider the hypercubic lattice $\Delta\mathbb{Z}^4$ with points $x=\Delta n$, $n\in\mathbb{Z}^4$.
A path from $0$ to $x$ is a sequence of $N$ steps $\nu_1,\dots,\nu_N$, each
$\nu_j\in\{\pm e_\mu\}_{\mu=1}^4$ (8 choices), with $\sum_j\nu_j=n$. Here $L(\text{path})=N$ is the
number of links, and revisiting links is allowed. Define
\begin{equation}
G(n)=\sum_{\text{paths }0\to x}e^{-m_0L(\text{path})}
=\sum_{N=0}^{\infty}e^{-m_0N}\,W_N(n),
\end{equation}
where $W_N(n)$ is the number of $N$-step walks from $0$ to $n$.

\section*{1. Counting paths with the Kronecker delta}

Use $\delta_{n,0}=\int_{-\pi}^{\pi}\frac{\dd k}{2\pi}e^{ikn}$ in each of the four directions to impose
the constraint $\sum_j\nu_j=n$:
\begin{equation}
W_N(n)=\sum_{\nu_1,\dots,\nu_N}\delta^{(4)}_{\sum_j\nu_j,\,n}
=\int_{[-\pi,\pi]^4}\frac{\dd^4k}{(2\pi)^4}\,e^{-ik\cdot n}
\prod_{j=1}^{N}\Big(\sum_{\nu_j}e^{ik\cdot\nu_j}\Big).
\end{equation}
Each factor is a sum over the 8 steps,
\begin{equation}
\sum_{\nu}e^{ik\cdot\nu}=\sum_{\mu=1}^{4}\big(e^{ik_\mu}+e^{-ik_\mu}\big)=2\sum_{\mu=1}^{4}\cos k_\mu ,
\end{equation}
so
\begin{equation}
W_N(n)=\int_{[-\pi,\pi]^4}\frac{\dd^4k}{(2\pi)^4}\,e^{-ik\cdot n}\Big(2\sum_{\mu=1}^4\cos k_\mu\Big)^{N}.
\end{equation}
(The integrand is symmetric under $k\to-k$, so the sign in the exponent is immaterial.)

\section*{2. Summing over $N$}

Summing the geometric series,
\begin{equation}
G(n)=\int_{[-\pi,\pi]^4}\frac{\dd^4k}{(2\pi)^4}\,
\frac{e^{ik\cdot n}}{1-2e^{-m_0}\sum_{\mu}\cos k_\mu}.
\label{eq:G}
\end{equation}
Since $\big|2\sum_\mu\cos k_\mu\big|\le 8$, the series converges for $8e^{-m_0}<1$, i.e.
$m_0>\ln 8$.

\paragraph{Behaviour at fixed $m_0$.}
\begin{itemize}
\item \textbf{$m_0<\ln 8$.} The number of walks grows like $8^N$ (e.g.\ $W_N(0)\sim 8^N N^{-2}$ for even $N$),
which beats $e^{-m_0N}$. The sum over $N$ diverges.
\item \textbf{$m_0>\ln 8$.} Every additional step is suppressed by a factor at most
$8e^{-m_0}<1$. The sum is dominated by the shortest paths, with $N_{\min}=\sum_\mu|n_\mu|=|x|_1/\Delta$, and
$G(n)$ decays exponentially, $G\sim e^{-\mu_{\rm lat}|n|}$ with $\mu_{\rm lat}>0$ fixed. As $\Delta\to0$ at fixed
$x$ we have $|n|\sim|x|/\Delta\to\infty$, so $G\to0$. In physical units the correlation length is
$\Delta/\mu_{\rm lat}\to0$.
\end{itemize}
For long paths to remain relevant, $m_0$ must approach the critical value $\ln 8$ as $\Delta\to0$,
so that the energy cost $m_0$ per link nearly cancels the entropy $\ln 8$ per link.

\section*{3. Tuning $m_0(\Delta)$}

Parametrize the distance to the critical value by a finite parameter $M$ (to be the physical mass):
\begin{equation}
e^{m_0}=8+M^2\Delta^2 .
\label{eq:tune}
\end{equation}
Using $8-2\sum_\mu\cos k_\mu=2\sum_\mu(1-\cos k_\mu)=4\sum_\mu\sin^2(k_\mu/2)$, the denominator in
\eqref{eq:G} is
\begin{equation}
1-2e^{-m_0}\sum_\mu\cos k_\mu
=e^{-m_0}\Big(M^2\Delta^2+4\sum_{\mu}\sin^2\tfrac{k_\mu}{2}\Big),
\end{equation}
so that, exactly,
\begin{equation}
G(n)=e^{m_0}\int_{[-\pi,\pi]^4}\frac{\dd^4k}{(2\pi)^4}\,
\frac{e^{ik\cdot n}}{M^2\Delta^2+4\sum_\mu\sin^2(k_\mu/2)} .
\end{equation}
This is the lattice Green's function of $-\nabla^2_{\rm lat}+M^2\Delta^2$.

Now rescale $k=\Delta p$, so that $e^{ik\cdot n}=e^{ip\cdot x}$, $\dd^4k=\Delta^4\dd^4p$ and
$p\in[-\pi/\Delta,\pi/\Delta]^4$:
\begin{equation}
G=e^{m_0}\Delta^{2}\int_{[-\pi/\Delta,\pi/\Delta]^4}\frac{\dd^4p}{(2\pi)^4}\,
\frac{e^{ip\cdot x}}{M^2+\hat p^{\,2}},\qquad
\hat p^{\,2}=\frac{4}{\Delta^2}\sum_\mu\sin^2\frac{\Delta p_\mu}{2}.
\end{equation}
For fixed $p$, $\hat p^{\,2}=p^2+O(\Delta^2p^4)$, and the Brillouin-zone boundary moves to infinity.
For $x\neq0$ the region $|p|\sim1/\Delta$ contributes negligibly because of rapid oscillations of
$e^{ip\cdot x}$. Since $e^{m_0}\to8$,
\begin{equation}
\boxed{\;G(n)\;\xrightarrow[\Delta\to0]{}\;8\Delta^2\,D_E(x;M),\qquad
D_E(x;M)=\int\frac{\dd^4p}{(2\pi)^4}\frac{e^{ip\cdot x}}{p^2+M^2}\;}
\end{equation}
which is the Euclidean Klein--Gordon correlation function. The factor $8\Delta^2$ is an overall
constant independent of $x$; it can be absorbed into the normalization of the (lattice) measure.

\section*{4. Result and interpretation}

Solving $e^{m_0}=8+M^2\Delta^2$ for $m_0$,
\begin{equation}
\boxed{\;m_0(\Delta)=\ln\!\big(8+M^2\Delta^2\big)=\ln 8+\frac{M^2\Delta^2}{8}+O(\Delta^4)\;}
\end{equation}
with $M$ held fixed as $\Delta\to0$. If instead $L$ is measured in physical length, $L=\Delta N$,
then $m_0\to m_0\Delta$ in the weight, i.e.\ $m_0=\frac{\ln 8}{\Delta}+\frac{M^2\Delta}{8}+\dots$

\paragraph{Proper-time check.}
Near $k=0$, $2e^{-m_0}\sum_\mu\cos k_\mu\simeq\big(1-\tfrac{M^2\Delta^2}{8}\big)\big(1-\tfrac{k^2}{8}\big)
\simeq \exp\!\big[-\tfrac{\Delta^2}{8}(p^2+M^2)\big]$. Hence the $N$-step contribution is weighted by
$e^{-s(p^2+M^2)}$ with proper time
\begin{equation}
s=\frac{N\Delta^2}{8},\qquad\text{i.e. } \xi=\tfrac18 ,
\end{equation}
and replacing $\sum_N\to\frac{8}{\Delta^2}\int_0^\infty\dd s$ gives
\begin{equation}
G\to 8\Delta^2\int_0^\infty\dd s\,e^{-sM^2}\int\frac{\dd^4p}{(2\pi)^4}e^{-sp^2+ip\cdot x}=8\Delta^2D_E(x;M),
\end{equation}
which is the Schwinger proper-time representation. This agrees with the result above and with the
tuning $\Delta m_0-\log A\to\xi m^2\Delta^2$ of Lecture~7 (the lattice analogue of $\log A$ is $\ln 8$).

Keeping $s$ finite requires $N\sim 8s/\Delta^2\to\infty$. The continuum limit is therefore not obtained
at fixed physical path length $N\Delta$: the dominant paths have infinitely many microscopic links
($N\sim 1/(M^2\Delta^2)$). They are very long, crumpled, Brownian-like paths (fractal dimension 2),
with typical extent $x^2_{\rm typ}\sim N\Delta^2\sim s$, and not the shortest path.

\subsection*{AI use declaration}

Model(s) used: Claude Sonnet 5.5 (Anthropic) and ChatGPT Luna 5.6 (OpenAI).\\

Revisions/corrections made to model output, and clarifications added:
The suggestions of the second model were incorporated: the wording about the overall factor
$8\Delta^2$ was made less assertive (it is an overall normalization constant, independent of $x$),
and the statement about the typical extent of the paths was changed to
$x^2_{\rm typ}\sim N\Delta^2\sim s$. In addition, I edited the \LaTeX{} source by hand and made some other small changes.\\

Comparisons/checks against the lecture notes:
The result was compared with the notes (the section on the path integral for a relativistic particle,
Lecture~7). The $\ln 8$ term (logarithm of the number of directions) plays the role of $\log A$,
the correction $M^2\Delta^2/8$ corresponds to $\xi m^2\Delta^2$ with $\xi=1/8$, and the continuum
limit reproduces the Schwinger proper-time representation of $D_E(x)$.\\

Calculation steps, limiting cases, or consistency checks verified by hand:
I checked by hand the following steps: the convergence condition $8e^{-m_0}<1$; the identity
$8-2\sum_\mu\cos k_\mu=4\sum_\mu\sin^2(k_\mu/2)$; the small-$k$ expansion
$2\sum_\mu\cos k_\mu\approx 8-k^2$; the rescaling $k=\Delta p$; and the proper-time identification
$s=N\Delta^2/8$.\\

Calculation steps, limiting cases, or consistency checks subjected to automated verification
(via Mathematica, python, LLMs):
The derivation was produced by Claude Sonnet 5.5 and reviewed by a second LLM (ChatGPT Luna 5.6).
No Mathematica or python verification was performed.

\end{document}